\documentclass{scrartcl}
\usepackage{amsmath}
\usepackage{natbib}
\usepackage{graphicx}
\usepackage{geometry}
\geometry{ a4paper, 
    left = 25mm,
    right = 25mm,
    top = 20mm,
    bottom = 25mm,
    footskip = 10mm,} 
\usepackage{setspace}
\onehalfspacing
\usepackage{tikz}
\usetikzlibrary{shapes}
\usepackage{hyperref}
\usepackage{booktabs}
\widowpenalty10000
\clubpenalty10000
\setlength\emergencystretch{3em}
\renewcommand{\textfraction}{0.1}
\renewcommand{\floatpagefraction}{0.8}
\renewcommand{\topfraction}{0.9} 
\renewcommand{\bottomfraction}{0.8} 
\renewcommand{\textfraction}{0.07} 
\renewcommand{\floatpagefraction}{0.7}  

\title{Måling af Variation i Profitabilitet  \\ -- i Dansk Planteproduktion}
\author{Helle Dahlgren Berlin og \\ William Louis Reingaard Vissing \\
    \small Institut for Fødevare- og Ressourceøkonomi, Københavns Universitet}

\date{25. februar, 2026}

\begin{document}

\maketitle

\noindent
Der er stigende konkurrence om arealanvendelsen i Danmark og et centralt spørgsmål er, hvordan det bedst sikres, at den mest produktive landbrugsjord forbliver i landbrugsmæssig anvendelse, samtidig med at reguleringsindsatser målrettes de mest miljøbelastende jorde. I dag udvikles der mange geografiske lag, der ud fra biologiske og miljømæssige hensyn angiver, hvor der eksempelvis bør plantes skov eller reduceres kvælstof – bl.a. i forbindelse med den grønne trepart. \\[1em]
Med henblik på at styrke beslutningsgrundlaget ønsker vi at udvikle en økonometrisk model, der estimerer dækningsbidraget på markniveau. Givet jordkvalitet, klimatiske forskelle mm., forventes der at være en betydelig geografisk variation i dyrkningsværdier for planteavl. Her påtænker vi at anvende den samme metode der bruges i produktionsgrensstatistikken, som IFRO – på vegne af Danmarks Statistik og Ministeriet for Fødevarer, Landbrug og Fiskeri – foretager hvert år med regnskabsdata fra danske landbrug til at allokere input mængder / omkostninger fra bedriftsniveau til produktionsgrensniveau. Metoden skal dermed tilpasses og anvendes med formålet om at allokere resultater fra bedriftsniveau til markniveau. Hertil har vi brug for adgang til udtræk fra SEGES økonomidatabase. \\[1em]
Derudover ønsker vi at inkludere kvælstofanvendelse og ikke mindst kvælstofudvaskning til kystvande på markniveau. Hertil har vi brug for SEGES bæredygtighedsdata særligt data for anvendelse af husdyrgødning. \\[1em]
Specialet forventes at fungere som et forstudie, der kan bidrage til et mere nuanceret beslutningsgrundlag for arealplanlægning og miljøregulering, herunder identificere marker med høj privatøkonomisk værdi og lav miljømæssig omkostning – og omvendt.



\end{document}